\documentclass[print]{styles/friggeri-cv-mac} % Add 'print' as an option into the square bracket to remove colors from this template for printing
%\documentclass[print]{styles/friggeri-cv-linux} % Add 'print' as an option into the square bracket to remove colors from this template for printing

\usepackage{soul}
\addbibresource{mylibrary.bib}


\begin{document}

\header{José Jair }{Santanna}{Enthusiastic, Teacher, Researcher, Network Security Specialist \& (Big) Data Analyst} 

\begin{aside} 
 \section{Contact}
University of Twente
DACS research group
Drienerlolaan 5 
7522 NB Enschede
The Netherlands
~
\href{mailto:j.j.santanna@utwente.nl}{j.j.santanna@utwente.nl}
+31 53 4892505
~
\section{More Information}\hspace{-1cm}
\href{http://www.jairsantanna.com}{jairsantanna.com}\includegraphics[scale=0.5]{img/jairsantanna.png}
\href{https://scholar.google.com/citations?user=TxcQNxUAAAAJ}{/TxcQNxUAAAAJ}\includegraphics[scale=0.3]{img/googlescholar.png}
\href{http://orcid.org/0000-0002-8361-6729}{/0...2-8361-6729}\includegraphics[scale=0.3]{img/orcid.png}
\href{http://www.researcherid.com/rid/K-2169-2016}{/K-2169-2016}\includegraphics[scale=0.3]{img/researchid.png}
\href{http://www.slideshare.net/jjcsantanna}{/jjcsantanna}\includegraphics[scale=0.3]{img/slideshare.png}
\href{https://www.linkedin.com/in/jjcsantanna}{/jjcsantanna}\includegraphics[scale=0.3]{img/linkedin.png}
\href{https://github.com/jjsantanna}{/jjsantanna\includegraphics[scale=0.3]{img/github.png}}
%\href{https://www.facebook.com/jjcsantanna}{/jjcsantanna}\includegraphics[scale=0.3]{img/facebook.png}
%~
%\section{Programming}
%	Bash, Python, C\#,SQL
%	PIG, Spark/Scala \& Impala
~
\section{Languages}
Dutch\includegraphics[scale=0.40]{img/2stars.png}
Spanish\includegraphics[scale=0.40]{img/3stars.png}
English\includegraphics[scale=0.40]{img/5stars.png}
Portuguese\includegraphics[scale=0.40]{img/5stars.png}
\end{aside}

\section{Professional Objective}\vspace{-10pt}
\setlength\parindent{12pt}I am currently a PhD candidate at University of
Twente with already a postdoctoral contract starting right after I achieve the
Doctor degree (march 2017). My research is founded by the European commission
and by the Dutch national Organisation for Scientific Research (NWO). Some
results of my PhD research has already achieved recognition in the national and
international level, for example I had the opportunity to advise both the
Japanese Ministry of Internal Affairs and Communications and the Dutch National
Cyber Security Center on the topic of cyber security. 

\vspace{-5pt}\setlength\parindent{12pt}My professional objective is to become a
recognized researcher and professor in network management and security. On top
of that I am eager to make my name and the name of the academic institution that
I am part of as a worldwide reference in cyber security. I aim to transfer my
enthusiasm, knowledge, skills and experience into meaningful innovation to
improve the security in all types of networks, specially in the Internet. 

\setlength\parindent{0pt}
\section{Education}\vspace{-5pt}

\begin{entrylist}

\entry
{2013--mar/2017*}
{Doctor {\normalfont in Computer Science}}
{University of Twente, The Netherlands}
{\emph{``The DDoS as a \$ervice Phenomenon''.} \\
This ongoing work have being used by both the security community and by network
operators. One of the outcomes that stands out is publicly available at
\href{http://booterblacklist.com}{`booterblacklist.com'} and was already downloaded by more
than four hundred users worldwide.\\ 
Supervisor: Prof. Dr. Ir. Aiko Pras }	

\entry
{2010--2012}
{Master {\normalfont in Computer Science}}
{Federal University of Rio Grande do Sul, Brazil}
{\emph{``A BPM-based Solution for Inter-domain Circuit Management''.}\\
The outcome of this work is still in production used by the Brazilian
National Education and Research Network (Link:
\href{http://meican.cipo.rnp.br/}{http://meican.cipo.rnp.br/}).\\ 
Supervisor: Prof. Dr. Lisandro Zambenedetti Granville}

\entry
{2005--2010}
{Bachelor {\normalfont in Computer Engineer}}
{Federal University of Par\'a, Brazil}
{\emph{``A Comparative Evaluation of DCCP Protocol in Mesh Networks.''} \\ 
Supervisor: Dr. Kelvin Lopes Dias
%This dissertation ...
}

\end{entrylist}

\section{Experience}\vspace{-5pt}

\subsection{Technical Adviser}\vspace{-5pt}
\begin{entrylist}
\entry
{2016}
{To the Yokohama National University}
{Japan}
{For the production of a report on ``Reflection \& Amplification DDoS attacks'' requested by the \emph{Japanese Ministry of Internal Affairs and Communications}.}

\entry
{2016}
{To the Dutch National Cyber Security Center and the Scientific Research \& Documentation Center}
{The Netherlands}
{For the production of a report on ``Cyber metrics 2016: DDoS and Malware'' requested by the \emph{Dutch Ministry of Security and Justice}.}
\end{entrylist}

\subsection{Assistant Professor}\vspace{-5pt}
\begin{entrylist}
\entry
{2012--2013}
{University of Santa Cruz do Sul (UNISC)}
{Santa Cruz do Sul, Brazil}
{Tasks \& Responsibilities: Supervise students and teach (integrally) the courses: \textit{Digital
Systems, Computer Networks, and Network Management}.} 
\end{entrylist}

\subsection{Guest Lecturer}\vspace{-5pt}
\begin{entrylist}
\entry
{2016}
{Network Management}
{University of Twente, The Netherlands}
{Title: SNMPv1, v2, v3 and Beyond (a very practical lecture)\\
Audience: master students of computer science, electrical engineer, and telematics from the University of Twente. }

\entry
{2016}
{Network Management}
{University of Twente, The Netherlands}
{Title: Network management based on WebServices\\
Audience: master students of computer science, electrical engineer, and telematics from the University of Twente. }

\entry
{2015 \& 2016}
{Product Design to Online Business}
{University of Twente, The Netherlands}
{Audience: bachelor students of industrial engineering and business information technology.}

\entry
{2015}
{Network Security}
{University of Twente, The Netherlands}
{Audience: master and bachelor students of computer science, electrical engineer, and telematics from the University of Twente, 3TU cyber security, and members of ICT labs. }

\entry
{2015}
{Cybercrime \& Cybersecurity (minor)}
{University of Twente, The Netherlands}
{Audience: bachelor students from University of Twente, European Research Center for Information (ERCIS), Westfälische Wilhelms - Univerität Münster, Universität Innsbruck, and University of Leicester. }
\end{entrylist}

\subsection{Student Supervisor}\vspace{-5pt}
Among 10 supervised students: 5 decided to follow towards the next academic degree[$^*$], PhD or MSc, 2 received the MSc degree cum laude[$^+$], 2 won the best paper award in international conferences[$^\#$], and 1 won the second place of a renowned Dutch national prize (KIVI/Defensie \& Veiligheid) for her thesis[$!$].

\begin{entrylist}
\vspace{-0.3cm}
\entry
{[M.Sc. Degree]}
{Justyna Chromik[$^*$ $^+$ $^\#$ $!$] , Wouter de Vries[$^*$ $^\#$], Joey de Vries[$^+$], Mark Wierbosch, and Jarmo van Lenthe.}
{}

\vspace{-0.3cm}
\entry
{[B.Sc. Degree]}
{Roeland Krak[$^*$], Max Kerkers[$^*$], Dirk Maan[$^*$], Jochum Börger, and Guilherme Dressler.}
{}

%[MINOR in Cyber Attacks and Cyber Security ] Karel Kroonen, Michiel Hamilton, Rick Clement, Peter Schroten, Yuri van Midden, Benedict Steinhoff, Kilian Müller, Stephanie Massa, Sören Schleibaum
%[DESIGN Project] Tom Velthausz, Henk Mulder, Thomas Reesink, Jordy Michorius,Dennis Eijkel;
%[Crossing Border Module]Julik Keijer

\end{entrylist}

\subsection{International Conference Reviewer}\vspace{-5pt}
Among more than 20 venues acting as reviewer we highlight four conferences which are the top well ranked related to the Committee on Network Operation and Management (CNOM), which is a renowned technical committee part of the IEEE Communications Society. For each conference we present the acronym, the full name, the acceptance rate which is the number of papers submitted divided by published* (according \href{http://www.cs.ucsb .edu/~almeroth/conf/stats}{\texttt{www.cs.ucsb.edu/$\textasciitilde$almeroth}}), and the h-index* (according Google scholar). For the entire list of conference that I was reviewer you can check my records on the known conference management systems, i.e., EDAS [link], JEMS [link], and EasyChair [link].

\begin{entrylist}

\vspace{-0.3cm}
\entry
{CNSM} %(2013-- 2015)
{International Conference on Network and Service Management}
{17.6\% | 18}

\vspace{-0.3cm}
\entry
{NETWORKING} %(2014--2016)
{IFIP Networking}
{23.8\% | 22 }

\vspace{-0.3cm}
\entry
{IM} %(2015)
{IFIP/IEEE Integrated Management}
{27.2\% | 18}

\vspace{-0.3cm}
\entry
{NOMS} %(2015)
{IEEE/IFIP Network Operations and Management Symposium }
{28.9\% | 22}


\end{entrylist}

\let\thefootnote\relax\footnotetext{*last checked: April 7,2016}

\newpage
\subsection{Conference Organizer}\vspace{-5pt}

\begin{entrylist}
\entry
{2016} 
{Shadow Technical Program Committee (TPC)}
{CNSM}
{IFIP/IEEE International Conference on Network and Service Management (CNSM) is the premier annual single-track conference in the general area of network and service management. A shadow TPC is a experimental initiative to make PhD researchers better TPCs in the future. The reviews are considered as normal a TPC review. }

\entry
{2013-2016} 
{Technical Program Committee (TPC)}
{ERRC}
{Escola Regional de Redes de Computadores, a conference for undergrad students from the south of Brazil.}

\entry
{2012} 
{Local Organizing Committee}
{LatinCloud}
{IEEE Latin American Conference on Cloud Computing and Communications}

\entry
{2010}
{Local Organizing Committee}
{SBRC}
{The Brazilian Symposium on Computer Networks and Distributed Systems.}

\end{entrylist}




%\subsection{Project Activities}\vspace{-5pt}
%Management of Hybrid Networks (HyMan)

%Desenvolvimento de Tecnologia para Redes Metropolitanas Sem Fio voltada para Serviços de Cidades Inteligentes, Descrição: O projeto tem como objetivo desenvolver uma solução integrada para os componentes de escalonamento e controle de admissão para dispositivos de 4 geração utilizando a tecnologia WiMAX.. , Situação: Em andamento; Natureza: Pesquisa. , Alunos envolvidos: Graduação: (2) . , Integrantes: José Jair Cardoso de Santanna - Integrante / Cristiano Bonato Both - Coordenador.

%2007 - 2009
%Arcabouço para Acesso ao Conhecimento à População da Região Amazônica, 
%Descrição: Abordamos aspectos de pesquisa nas áreas de infra-estrutura de comunicações, implementação e avaliação de interfaces amigáveis para acesso universal aos cidadãos, considerando as oportunidades de conectividade existentes e peculiaridades da região amazônica. O estudo e solução propostos neste projeto consideram hardware de comunicação de baixo custo, além de software livre e aberto para a composição de um sistema de acesso flexível e escalável.. , Situação: Concluído; Natureza: Pesquisa. , Alunos envolvidos: Graduação: (3) / Mestrado acadêmico: (7) / Doutorado: (2) . , Integrantes: José Jair Cardoso de Santanna - Integrante / Kelvin Lopes Dias - Coordenador., Financiador(es): Conselho Nacional de Desenvolvimento Científico e Tecnológico - Auxílio financeiro.


\subsection{Software Developer}\vspace{-5pt}
\begin{entrylist}

\entry
{2015--Now}
{Distributed Denial of Service Attacks Database (DDoSDB)}
{\\Link: \href{http://ddosdb.org}{http://www.ddosdb.org} }
{An initiative to develop the largest open database with DDoS attacks.\\
Languages: Python, Bash, Pig Latin \& Apache Impala
}

\entry
{2010--2011}
{{Management Environment of Inter-domain Circuits for Advanced Networks (MEICAN)}}
{Link: \href{http://meican.cipo.rnp.br/}{http://meican.cipo.rnp.br/}}
{Development of an Executable Network Management Policies Editor sponsored by the Brazilian National Education and Research Network (RNP).    \\
Language: WS-BPEL
}

\entry
{2009}
{Support for Planning Sustainable Electrical Power Generation via Vegetable Oil.}
{}
{Desktop application part of a Ph.D. thesis [\href{https://sites.google.com/site/ceamazonufpa/TeseAnaRosaDuarte.pdf}{link}] sponsored by a subsidiary (Eletronorte) of the major Brazilian power utility (Eletrobrás).\\
Language: C\#.Net \& mySQL
}
\end{entrylist}

\section{Awards}\vspace{-5pt}

\begin{entrylist}
\vspace{-0.3cm}
\entry
{2015}
{Best Ph.D. 'Elevator Pitch' (1$^{st}$ Place)}
{CTIT Symposium, University of Twente}

\vspace{-0.3cm}
\entry
{2015}
{Best Poster (3$^{rd}$ Place)}
{CTIT Symposium, University of Twente}

\vspace{-0.3cm}
\entry
{2015}
{Best Paper Award}
{Autonomous Infrastructure, Management and Security Conference}

\vspace{-0.3cm}
\entry
{2015}
{Top Papers Nomination}
{Brazilian Symposium on Computer Networks and Distributed Systems}

\vspace{-0.3cm}
\entry
{2010}
{Student Award (3$^{rd}$ Place)}
{Computer Engineer, Federal University of Pará, Brazil}

\end{entrylist}

\newpage
\section{Presentations}\vspace{-5pt}
%most reviewers will not take time to check this website, so you have to give a bit more information. For instance how many at conferences, which conferences, how many invited, where? etc. 
\emph{More than 30 presentations in the last 3 years (available on my  \href{http://www.slideshare.net/jjcsantanna}{slideshare.net/jjcsantanna}), among those I highlight:}

\begin{entrylist}

\vspace{-0.3cm}
\entry
{2015}
{Invited presentation}
{at Logius (part of the Dutch Ministry of the Interior and Kingdom Relations).}

\vspace{-0.3cm}
\entry
{2015}
{Invited Presentation}
{to the Dutch National Management Organization of Internet Providers (NBIP).}

\vspace{-0.3cm}
\entry
{2015}
{Invited Presentation}
{to the Dutch National Cyber Security Centrum (NCSC).}

\vspace{-0.3cm}
\entry
{2013 \& 2015}
{Presentation}
{at the 31$^{th}$ and 36$^{th}$ IRTF Network Management Research Group (NMRG).}

\end{entrylist}

\section{Capacity to Stablish Collaboration}\vspace{-5pt}
\url{Booking.com}, \url{Nationale Beheersorganisatie Internet Providers(NBIP)},
\url{fox-it.com}, \url{quarantainenet}, \url{SURFNet}, \url{CESNET}, \url{RIPE NCC}




\let\thefootnote\relax\footnotetext{CV last update: \today}

\newpage
\section{Scientific Appendix}\vspace{-5pt}

\subsection{Publications}

\nocite{*}
\printbibliography[title={\emptyset}]

\end{document}
